\documentclass{hcmut-report}
\usepackage{subfiles}
\usepackage{codespace}
\usepackage{float}
\usepackage{tikz}
\usetikzlibrary{positioning}
\usepackage{gensymb,textcomp}
\usepackage{array,longtable,multicol,multirow,tabularx}
\usepackage{enumitem}
\usepackage{caption}
\usepackage[export]{adjustbox}
\usepackage[ruled,vlined,linesnumbered]{algorithm2e}

\setlength{\emergencystretch}{3em}

\SetAlgorithmName{Thuật toán}{thuật toán}{Danh sách thuật toán}
\SetKwInput{KwIn}{Đầu vào}
\SetKwInput{KwOut}{Đầu ra}

% Thông tin trang bìa
\coursename{BÁO CÁO BÀI TẬP LỚN\\ MÔN KHAI PHÁ DỮ LIỆU (CO3029)}
\title{PHÂN ĐOẠN KHÁCH HÀNG ĐA CHIỀU THEO MÔ HÌNH RFM KẾT HỢP HÀNH VI}
\major{Khoa học Máy tính}
\supervisorone{Bùi Tiến Đức}
\studentone{Bùi Đình Phúc -- 2312664 (nhóm trưởng)}
\studenttwo{Hoàng Hữu Nhân -- 2312419}
\studentthree{Nguyễn Cảnh Nguyên -- 2312357}
\studentfour{Trương Hoàng Danh -- 2310423}
\reportplace{TP. Hồ Chí Minh}
\reporttime{tháng 11 năm 2026}

\AtBeginDocument{\counterwithin{lstlisting}{section}}

\begin{document}
\coverpage%
\begingroup
\setstretch{1}
\tableofcontents
\endgroup
\clearpage

\subfile{chapters/01_gioi_thieu.tex}
\subfile{chapters/02_du_lieu_eda.tex}
\subfile{chapters/03_tien_xu_ly.tex}
\subfile{chapters/04_mo_hinh.tex}
\subfile{chapters/05_danh_gia.tex}
\subfile{chapters/06_ung_dung.tex}
\subfile{chapters/07_ket_luan.tex}

\clearpage
\phantomsection
\addcontentsline{toc}{section}{TÀI LIỆU THAM KHẢO}
\renewcommand{\refname}{TÀI LIỆU THAM KHẢO}
\begin{thebibliography}{99}
\bibitem{chen2019dataset} D. Chen. \emph{Online Retail II} [Dataset]. UCI Machine Learning Repository, 2019. \url{https://doi.org/10.24432/C5CG6D}.
\bibitem{chen2012rfm} D. Chen, S. L. Sain, K. Guo. Data mining for the online retail industry: A case study of RFM model-based customer segmentation using data mining. \emph{Journal of Database Marketing \& Customer Strategy Management}, 19(3):197--208, 2012.
\bibitem{hughes1994} A. M. Hughes. \emph{Strategic Database Marketing}. Probus Publishing, 1994.
\bibitem{han2022} J. Han, J. Pei, H. Tong. \emph{Data Mining: Concepts and Techniques}, 4th ed. Morgan Kaufmann, 2022.
\bibitem{lloyd1982} S. Lloyd. Least squares quantization in PCM. \emph{IEEE Transactions on Information Theory}, 28(2):129--137, 1982.
\bibitem{arthur2007} D. Arthur, S. Vassilvitskii. k-means++: The advantages of careful seeding. In \emph{Proceedings of the 18th ACM-SIAM Symposium on Discrete Algorithms}, pp. 1027--1035, 2007.
\bibitem{ward1963} J. H. Ward. Hierarchical grouping to optimize an objective function. \emph{Journal of the American Statistical Association}, 58(301):236--244, 1963.
\bibitem{ester1996} M. Ester, H.-P. Kriegel, J. Sander, X. Xu. A density-based algorithm for discovering clusters in large spatial databases with noise. In \emph{Proceedings of KDD-96}, pp. 226--231, 1996.
\bibitem{rousseeuw1987} P. J. Rousseeuw. Silhouettes: A graphical aid to the interpretation and validation of cluster analysis. \emph{Journal of Computational and Applied Mathematics}, 20:53--65, 1987.
\bibitem{davies1979} D. L. Davies, D. W. Bouldin. A cluster separation measure. \emph{IEEE Transactions on Pattern Analysis and Machine Intelligence}, 1(2):224--227, 1979.
\bibitem{calinski1974} T. Cali\'{n}ski, J. Harabasz. A dendrite method for cluster analysis. \emph{Communications in Statistics}, 3(1):1--27, 1974.
\bibitem{hubert1985} L. Hubert, P. Arabie. Comparing partitions. \emph{Journal of Classification}, 2(1):193--218, 1985.
\bibitem{pedregosa2011} F. Pedregosa et al. Scikit-learn: Machine learning in Python. \emph{Journal of Machine Learning Research}, 12:2825--2830, 2011.
\end{thebibliography}

\end{document}
